% ============================================================
%  Chapter: Command-Line Interfaces, Virtual Environments,
%           and Shell Scripting for Windows and Linux
% ============================================================
%  Dependencies (add to your preamble if not already present):
%    \usepackage{listings}
%    \usepackage{xcolor}
%    \usepackage{booktabs}
%    \usepackage{hyperref}
%    \usepackage{tcolorbox}
%    \usepackage{fontawesome5}   % optional, for icons
% ============================================================


% ============================================================
\chapter{Command-Line Interfaces, Virtual Environments,
         and Shell Scripting}
\label{ch:cli-scripting}
% ============================================================

% ============================================================
\section{What Is a File?}
\label{sec:what-is-a-file}
% ============================================================

\subsection{Files and text}

A file is a named sequence of bytes stored on a storage device. From the
computer's perspective, all files are sequences of bytes. What distinguishes
a Python script from a photograph is not the storage mechanism but the
interpretation: a text editor interprets bytes as characters; an image
viewer interprets them as pixel values. A Python file is plain text: every
character you type is stored as a byte, and the Python interpreter reads
those bytes as instructions. This is why a Word document cannot serve as a
Python program; Word stores invisible formatting bytes alongside the text.

\subsection{The filesystem: folders and paths}

Every file lives inside a folder (also called a directory). Folders can
contain other folders, forming a tree. The location of any file can be
described by a \emph{path}: a sequence of folder names separated by a
delimiter, ending in the filename.

\begin{table}[htbp]
\centering
\caption{Path syntax on each operating system.}
\label{tab:path-syntax}
\begin{tabular}{p{3.5cm} p{9cm}}
\toprule
\textbf{Operating System} & \textbf{Example path} \\
\midrule
Windows & \texttt{C:\textbackslash Users\textbackslash Alice\textbackslash Documents\textbackslash hello.py} \\
macOS   & \texttt{/Users/alice/Documents/hello.py} \\
Linux   & \texttt{/home/alice/Documents/hello.py} \\
\bottomrule
\end{tabular}
\end{table}

Windows uses backslashes and a drive letter; macOS and Linux use forward
slashes and start from a root \texttt{/}.

\subsection{The working directory}

At any moment, the CLI is ``inside'' a particular folder, called the
\emph{current working directory} (CWD). Commands that refer to files
without a full path are interpreted relative to that folder. When you type
\texttt{python hello.py}, the interpreter looks for \texttt{hello.py} in
the CWD. If the file is elsewhere, the command fails. Understanding the
CWD is the single most common source of confusion for new CLI users and
the root cause of most ``file not found'' errors.

The CLI provides the commands to navigate the filesystem, inspect its
contents, and execute files; these are introduced in the next section.

% ============================================================

\begin{quote}
\textit{``The command line is the most powerful interface ever invented.
It is also the most intimidating. This chapter aims to change that.''}
\end{quote}

% ============================================================
\section{Introduction: What Is a Command-Line Interface?}
\label{sec:cli-intro}
% ============================================================

When you use a computer, you almost certainly interact with it through a
\textbf{Graphical User Interface} (GUI): windows, icons, menus, and
buttons that you click with a mouse. GUIs are intuitive and visually rich,
but they have a fundamental limitation: every action you can take must first
be anticipated and implemented as a button or menu item by the software
developer. If you need to do something the GUI was not designed for, you are
stuck.

A \textbf{Command-Line Interface} (CLI) solves this by letting you
communicate with the computer in \emph{text}. You type a command; the
computer executes it and responds in text. This is not a step backward; it
is a different, and in many contexts more powerful, mode of interaction.

\subsection{A Concrete Analogy}

Think of the difference between a restaurant with a fixed menu and a
restaurant where you can describe exactly what you want to the chef. A GUI
is the fixed menu: fast, friendly, and good enough for most customers. A CLI
is talking directly to the chef: it requires more knowledge of ingredients,
but the results can be precisely what you need, and you can combine them in
ways the menu never anticipated.

\subsection{Why Developers and Researchers Use the CLI}

\begin{itemize}
  \item \textbf{Automation.} A sequence of commands can be saved in a file
        (called a \emph{script}) and replayed perfectly every time.
  \item \textbf{Precision.} Every option and parameter is explicit and
        readable.
  \item \textbf{Remote access.} Servers in the cloud rarely have a graphical
        desktop; the CLI is often the only way in.
  \item \textbf{Speed.} Renaming 10{,}000 files takes one command; doing it
        by hand in a GUI might take hours.
  \item \textbf{Reproducibility.} You can share a script with a colleague
        and know they will get exactly the same result.
\end{itemize}

\subsection{The Two Major CLI Worlds}

This chapter focuses on two CLI environments:

\begin{description}
  \item[Windows Command Prompt / PowerShell] The native CLI tools in
        Microsoft Windows. Batch files (\texttt{.bat}) and PowerShell
        scripts (\texttt{.ps1}) automate tasks here.
  \item[Bash (Linux / macOS)] The dominant CLI on Linux systems and macOS.
        Shell scripts (\texttt{.sh}) automate tasks here.
\end{description}

These two worlds use different syntax, different path conventions
(backslash \texttt{\textbackslash} vs.\ forward slash \texttt{/}), and
different scripting philosophies. A key goal of this chapter is to help you
move fluidly between them.

\subsection{Your First Commands}

The best way to demystify the CLI is to try it. Table~\ref{tab:first-commands}
shows equivalent commands in both environments.

\begin{table}[h]
\centering
\caption{Equivalent CLI commands: Windows vs.\ Linux/macOS}
\label{tab:first-commands}
\begin{tabular}{@{}lll@{}}
\toprule
\textbf{Task} & \textbf{Windows (CMD)} & \textbf{Linux / macOS (Bash)} \\
\midrule
Print current directory  & \texttt{cd}              & \texttt{pwd}             \\
List files               & \texttt{dir}             & \texttt{ls -la}          \\
Change directory         & \texttt{cd foldername}   & \texttt{cd foldername}   \\
Go up one level          & \texttt{cd ..}           & \texttt{cd ..}           \\
Create a folder          & \texttt{mkdir myfolder}  & \texttt{mkdir myfolder}  \\
Delete a file            & \texttt{del file.txt}    & \texttt{rm file.txt}     \\
Copy a file              & \texttt{copy a.txt b.txt}& \texttt{cp a.txt b.txt}  \\
Print text to screen     & \texttt{echo Hello}      & \texttt{echo Hello}      \\
Clear the screen         & \texttt{cls}             & \texttt{clear}           \\
\bottomrule
\end{tabular}
\end{table}

\begin{tipbox}
In both environments, pressing the \textbf{Up arrow} key recalls your
previous commands, and pressing \textbf{Tab} auto-completes file and
directory names. These two habits alone will dramatically speed up your
work at the CLI.
\end{tipbox}

% ============================================================
\section{Two Worlds: Windows and Linux/Mac}
\label{sec:two-worlds}
% ============================================================

Before writing any scripts, it helps to understand the philosophical
differences between how Windows and Linux handle the command line. These
differences are not arbitrary; they reflect decades of divergent design
history.

\subsection{Paths and Separators}

Windows uses \emph{backslashes} and drive letters to describe file
locations:
\begin{lstlisting}[style=windowsStyle,escapeinside={(*@}{@*)}]
C:\Users\DrG\Documents\(*@\ProjectName@*)\main.py
\end{lstlisting}

Linux uses \emph{forward slashes} and a single unified file tree rooted
at \texttt{/}:
\begin{lstlisting}[style=linuxStyle,escapeinside={(*@}{@*)}]
/home/drg/(*@\ProjectName@*)/main.py
\end{lstlisting}

\subsection{Scripts and Executability}

On Windows, the file extension determines whether a file is executable.
Files ending in \texttt{.bat}, \texttt{.cmd}, or \texttt{.exe} are treated
as programs.

On Linux, the extension is irrelevant. Executability is a \emph{permission}
stored on the file itself. A file must be explicitly granted execute
permission before the system will run it. This is done with the
\texttt{chmod} command, which we cover in Section~\ref{sec:chmod}.

\subsection{Environment Variables and PATH}

Both systems use an environment variable called \texttt{PATH} to find
programs. When you type a command, the shell searches each directory listed
in \texttt{PATH} in order and runs the first match it finds.

\begin{itemize}
  \item \textbf{Windows:} PATH is managed through System Properties $\to$
        Environment Variables, or via the \texttt{setx} command.
  \item \textbf{Linux:} PATH is set in configuration files such as
        \texttt{\textasciitilde/.bashrc} (for the Bash shell) or
        \texttt{\textasciitilde/.zshrc} (for Zsh), which are read every time
        a new terminal session starts.
\end{itemize}

% ============================================================
\section{WSL: A Linux Environment Inside Windows}
\label{sec:wsl}
% ============================================================

For Windows users who want the power of Linux scripting without leaving their
familiar environment, Microsoft provides the
\textbf{Windows Subsystem for Linux} (WSL). WSL is not an emulator or a
virtual machine in the traditional sense; it is a compatibility layer built
into Windows that runs a real Linux kernel alongside the Windows kernel.

\subsection{What WSL Gives You}

\begin{itemize}
  \item A full Bash shell (or Zsh, Fish, etc.)
  \item A complete Ubuntu (or other distribution) file system
  \item Access to thousands of Linux command-line tools via \texttt{apt}
  \item The ability to read and write Windows files from Linux
        (your Windows \texttt{C:} drive appears at \texttt{/mnt/c/})
  \item GPU access for machine learning workloads
\end{itemize}

\subsection{Installing WSL}

Open PowerShell or CMD \emph{as Administrator} and run:

\begin{lstlisting}[style=windowsStyle, caption={Installing WSL from PowerShell}]
wsl --install
\end{lstlisting}

This installs WSL 2 and the Ubuntu distribution by default. After a
restart, a Ubuntu terminal will open and ask you to create a username and
password. That account is independent of your Windows login.

\begin{notebox}
Throughout the remainder of this chapter, all Linux instructions apply
equally to WSL on Windows and to a native Linux machine. Treat them as the
same environment.
\end{notebox}

\subsection{The Elegant Solution: Visual Studio Code and the WSL Extension}
\label{subsec:vscode-wsl}

One of the most powerful aspects of WSL is its integration with
\textbf{Visual Studio Code} (VS Code), Microsoft's free code editor.
Installing the \textbf{WSL extension} in VS Code creates a seamless bridge
between your Windows desktop and your Linux environment:

\begin{itemize}
  \item You open VS Code from within the WSL terminal by typing \texttt{code .}
  \item VS Code runs its interface on Windows but executes all code,
        terminal commands, extensions, and debuggers \emph{inside Linux}
  \item The file explorer shows your Linux file system
  \item The integrated terminal is a Bash shell
  \item Python environments, linters, and formatters all resolve to their
        Linux versions
\end{itemize}

This means you never have to choose between the Linux scripting ecosystem
and the Windows desktop experience: you get both simultaneously. From this
point forward, when we say ``open VS Code,'' we mean launching it from
the WSL terminal with \texttt{code .}

\begin{tipbox}
To install the WSL extension, open VS Code, press
\texttt{Ctrl+Shift+X} to open the Extensions panel, and search for
``WSL'' (publisher: Microsoft). One click installs it.
\end{tipbox}

% ============================================================
% End of chapter
% ============================================================